\section*{Datos y Fórmulas}

\textbf{Datos:}
\begin{equation}
    \lambda_1 = 33 \, \text{R/h}
\end{equation}
\begin{equation}
    d_1 = 1 \, \text{m}
\end{equation}

\textbf{Fórmula:}
\begin{equation}
    \lambda_2 = \left(\frac{d_1}{d_2}\right)^2 \lambda_1
\end{equation}
\begin{equation}
    \lambda_2 = \left(\frac{100}{10}\right)^2 (33) = 3,300 \, \text{R/h}
\end{equation}

Si tardo 10 segundos en meter la fuente, ¿cuál será la dosis que recibo en la mano?

Si recibiera 3300 R en una hora:
\begin{equation}
    \text{Dosis} = \frac{3300 \, \text{R}}{3600 \, \text{s}} \cdot 10 \, \text{s} = 9.17 \, \text{R}
\end{equation}

\subsection*{Cálculo de radiación de una fuente de Ir-192}
Una fuente de Ir-192 que fue comprada hace 148 días con 120 Ci, y está cubierta con 8" de Plomo, ¿qué nivel de radiación nos dará a 2 m?

\textbf{Primero calculamos la actividad de la fuente:}
- 148 días corresponden a dos vidas medias.
\begin{itemize}
    \item[$\circ$] 120 Ci inicialmente.
    \item[$\circ$] 60 Ci a los 74 días (primera vida media).
    \item[$\circ$] 30 Ci a los 148 días (segunda vida media).
\end{itemize}

\textbf{Ahora calculamos el nivel de exposición a 1 metro, sin blindaje.}
\begin{equation}
    30 \, \text{Ci} \times 0.55 = 16.50 \, \text{R/h} \, \text{a 1 m sin blindaje}
\end{equation}
